\chapter{Background and Related Work}\label{ch:literature}
% LENGTH: 7-10 pages.
% SOURCE: notes/related.md is already written in prose. Adapt it, do not
% rewrite from scratch. Plan.md section 4 has the reading list.
% RULE: published numbers from other papers come from different datasets and
% protocols. Quote them as context only, never inside our results tables.

\section{Image Restoration with CNNs}\label{sec:lit-cnn}
% DnCNN (residual learning) and FFDNet (noise-level map, pixel unshuffle).
% Mention NAFNet as the efficiency reference and Restormer/SwinIR as the
% heavy end. Keep these two short.
\todo{adapt from notes/related.md "Foundations"}

\section{Frequency-Domain Restoration}\label{sec:lit-freq}
% MWCNN (wavelet subbands), FFC, SFNet/FSNet (closest neighbour, spend the
% most words here), FFTformer. Spectral bias (Rahaman) for motivation.
% The positioning sentence must be precise: SFNet does learned local frequency
% selection; we use FIXED radial masks then learned spatial expert mixing.
\todo{adapt}

\section{Mixture of Experts and Gating}\label{sec:lit-moe}
% Shazeer sparse MoE, Switch Transformer routing stability, load balancing
% vs our entropy regulariser. SE networks for the fusion block.
% Be honest: we compute every expert, so we get sparse WEIGHTS, not sparse
% COMPUTE. Say it here so the examiner does not find it later.
\todo{adapt}

\section{Evaluation Metrics}\label{sec:lit-metrics}
% PSNR, SSIM, LPIPS: one paragraph each with the formula for PSNR and what
% LPIPS measures, and why it can disagree with PSNR.
\todo{write}

\section{Research Gap}\label{sec:gap}
% Comparison table: rows = methods, cols = frequency decomposition? learned
% routing? spatial routing? params. Keep the old proposal gap table OUT unless
% MWCNN/SFNet/FFTformer are added to it (plan.md section 0.1).
\begin{table}[htbp]
	\centering
	\caption{Positioning of \fcsg{} against closely related methods.}
	\label{tab:gap}
	\begin{tabular}{@{}lcccl@{}}
		\toprule
		Method & Freq.\ split & Learned routing & Spatial & Notes \\
		\midrule
		DnCNN \cite{zhang2017dncnn}     & -- & -- & -- & residual CNN \\
		FFDNet \cite{zhang2018ffdnet}   & -- & -- & -- & noise map input \\
		MWCNN \cite{liu2018mwcnn}       & wavelet & -- & -- & joint subbands \\
		SFNet \cite{cui2023sfnet}       & learned & \todo{check} & \todo{check} & closest prior \\
		Sparse MoE \cite{shazeer2017moe}& -- & yes & -- & token routing \\
		\fcsg{} (ours)                  & radial FFT & top-2 & per pixel & 97k params \\
		\bottomrule
	\end{tabular}
\end{table}
